Tahap pertama dalam penelitian ini adalah studi literatur untuk memahami konsep dasar dan teknologi yang relevan dengan topik penelitian.
Studi literatur mencakup kajian terhadap penelitian terdahulu mengenai segmentasi visual, fusi sensor dengan Kalman Filter, navigasi berbasis waypoint, serta metode perencanaan lintasan untuk kendaraan otonom.

Setelah memahami landasan teori, tahap berikutnya adalah pengumpulan dan pengolahan data.
Data segmen jalan dibuat menggunakan Google My Maps dan diekspor dalam format KML yang kemudian dikonversi menjadi waypoint dalam format GeoJSON.
Data visual dikumpulkan melalui perekaman kamera depan kendaraan pada berbagai kondisi pencahayaan dan struktur jalan di lingkungan kampus ITS.
Secara paralel, data sensor navigasi (GPS, IMU, dan odometri) direkam untuk keperluan pengembangan dan pengujian modul lokalisasi.

Tahap perancangan sistem meliputi empat subsistem utama.
Pertama, subsistem persepsi visual yang mencakup pelatihan model InternImage dengan \emph{fine-tuning} pada \emph{dataset} lokal ITS, konversi model ke \emph{engine} TensorRT, dan implementasi modul konversi masker biner ke \emph{point cloud}.
Kedua, modul lokalisasi menggunakan Kalman Filter 4-keadaan yang memfusikan data GPS, IMU, dan odometri roda.
Ketiga, perencanaan global berbasis graf segmen jalan dan algoritma Dijkstra.
Keempat, perencanaan lokal berbasis \emph{Dynamic Window Approach} dengan \emph{arc fan} yang mengintegrasikan data visual dari \emph{point cloud} dan rute global GPS.

Tahap implementasi mencakup integrasi seluruh modul perangkat lunak dalam kerangka ROS2, pemasangan sensor dan komputer pada kendaraan Omoda E5, serta konfigurasi komunikasi CAN FD untuk kendali kendaraan.
Kalibrasi kamera menggunakan metode Scaramuzza dilakukan untuk memperoleh parameter proyeksi BEV.

Tahap pengujian dilakukan dalam empat kategori.
Pengujian subsistem persepsi visual mengevaluasi akurasi segmentasi dan laju inferensi.
Pengujian subsistem lokalisasi mengevaluasi akurasi posisi terhadap data aktual dan terhadap DLIO pada berbagai kondisi jalan.
Pengujian subsistem perencanaan lokal mengevaluasi kemampuan menghasilkan lintasan pada berbagai geometri jalan.
Pengujian integrasi sistem mengevaluasi keberhasilan penyelesaian rute, akurasi destinasi, dan ketahanan terhadap gangguan GPS.

Tahap terakhir adalah analisis hasil dan penyusunan laporan tugas akhir yang mencakup seluruh proses penelitian dari perancangan hingga pengujian.

Rencana jadwal kegiatan pengerjaan tugas akhir dapat dilihat pada Tabel~\ref{tbl:timeline} berikut:

\newcommand{\w}{}
\newcommand{\G}{\cellcolor{gray}}
\begin{table}[H]
  \caption{Timeline Kegiatan Penelitian}
  \label{tbl:timeline}
  \centering
  \resizebox{\textwidth}{!}{%
  \begin{tabular}{|p{5.5cm}|c|c|c|c|c|c|c|c|c|c|c|c|c|c|c|c|}
    \hline
    \multirow{2}{*}{Kegiatan} & \multicolumn{16}{c|}{Minggu} \\
    \cline{2-17}
    & 1 & 2 & 3 & 4 & 5 & 6 & 7 & 8 & 9 & 10 & 11 & 12 & 13 & 14 & 15 & 16 \\
    \hline

    Studi Literatur &
    \G & \G & \G & \G & \w & \w & \w & \w & \w & \w & \w & \w & \w & \w & \w & \w \\
    \hline

    Pengumpulan Data &
    \w & \w & \G & \G & \G & \G & \w & \w & \w & \w & \w & \w & \w & \w & \w & \w \\
    \hline

    Perancangan Sistem &
    \w & \w & \w & \G & \G & \G & \G & \w & \w & \w & \w & \w & \w & \w & \w & \w \\
    \hline

    Implementasi &
    \w & \w & \w & \w & \w & \w & \G & \G & \G & \G & \w & \w & \w & \w & \w & \w \\
    \hline

    Pengujian &
    \w & \w & \w & \w & \w & \w & \w & \G & \G & \G & \G & \G & \w & \w & \w & \w \\
    \hline

    Analisis Hasil &
    \w & \w & \w & \w & \w & \w & \w & \w & \w & \w & \G & \G & \G & \G & \w & \w \\
    \hline

    Penulisan Laporan &
    \w & \w & \w & \w & \w & \w & \w & \w & \w & \w & \w & \w & \G & \G & \G & \G \\
    \hline
  \end{tabular}
  } % akhir resizebox
\end{table}
